% day3_seg2_inertial/straperr.tex
\begin{frame}
  \onslide<1->{
	\frametitle{Inertial Error Dynamics} 
	\framesubtitle{From Strapdown Equations to Error Equations}
}
\begin{itemize}
  \onslide<2->{
 \item Nav equations require perfect inertial instruments
}
  \onslide<3->{
 \item Instruments are not perfect
}
  \onslide<4->{
 \item Inertial instrument error sources (bias, scale factor errors, etc.) generate erroneous rate and acceleration inputs to the nav equations  
}
  \onslide<5->{
 \item Need to figure out the effects on the outputs of these equations
}
\end{itemize}
\end{frame}
% cm_kf_talk_straperr.tex


\begin{frame}
  \onslide<1->{
	\frametitle{Inertial Error Dynamics} 
	\framesubtitle{Strapdown Error Equations}
}
\begin{itemize}
  \onslide<2->{
 \item Nature of Errors
}
  \onslide<3->{
 \item Error Operator
}
  \onslide<4->{
 \item Error Equations
}
\end{itemize}
\end{frame}


\begin{frame}
  \onslide<1->{
	\frametitle{Inertial Error Dynamics} 
	\framesubtitle{Strapdown Error Equations: Nature of Errors}
}
\begin{itemize}
  \onslide<2->{
 \item Errors propagate according to derived equations
}
  \onslide<3->{
 \item Errors are generally small, so approximating assumptions
 help simplify the math
 \begin{itemize}
}
  \onslide<4->{
 \item linearize otherwise non-linear dynamics
}
  \onslide<5->{
 \item small angle assumptions:  $\sin\theta \approx \theta$
 \end{itemize}
}
\end{itemize}
\end{frame}


\begin{frame}
  \onslide<1->{
	\frametitle{Inertial Error Dynamics} 
	\framesubtitle{Strapdown Error Equations: Error Operator}
}
\begin{itemize}
  \onslide<2->{
\item   If $A$ is any quantity, we define error operator as follows:
\begin{equation}\label{errop}
  \errop{A},
\end{equation}
where $A$ is the true quantity of interest and $\hat{A}$ is an
apparent quantity.  The latter could be a measured value, and
estimated value, or some computed value.  This means that
\[\hat{A}=A+\delta A\]
i.e., the apparent value is the true value plus some error.

}
\end{itemize}
\end{frame}


\begin{frame}
  \onslide<1->{
	\frametitle{Inertial Error Dynamics} 
	\framesubtitle{Strapdown Error Equations: Error Operator}
}
\begin{itemize}
  \onslide<2->{
\item   Error operator satisfies:
\begin{equation}\label{erropid}
  \delta(AB) = \delta A \cdot B + A \cdot \delta B
\end{equation}
}
  \onslide<3->{
\item   Proof:
\begin{align*}
\delta(AB) &= \widehat{AB} - AB\\
&= \hat{A}\hat{B}-AB\\
&= (A+\delta A)(B+\delta B)-AB\\
&= AB + A\cdot\delta B +\delta A \cdot B + \delta A \cdot \delta B - AB\\
& = \delta A \cdot B + A\cdot\delta B
\end{align*}
since we may assume that $\delta A \cdot\delta B \approx 0$.
}
\end{itemize}
\end{frame}


\begin{frame}
  \onslide<1->{
	\frametitle{Inertial Error Dynamics} 
	\framesubtitle{Strapdown Error Equations}
}
\begin{itemize}
  \onslide<2->{
\item   Applying these to the three differential equations, yields three
error equations:
\begin{equation}\label{cblerreq}
  \cblerreq
\end{equation}
\begin{equation}\label{strapvelerreq}
  \strapvelerreq
\end{equation}
\begin{equation}\label{cleerreq}
  \cleerreq
\end{equation}
}
\end{itemize}
\end{frame}



\begin{frame}
  \onslide<1->{
	\frametitle{Inertial Error Dynamics} 
	\framesubtitle{Strapdown Error Equations}
}
\begin{itemize}
  \onslide<2->{
\item   Introducing \textbf{small angular rotations} for tilt and position
\begin{itemize}
}
  \onslide<3->{
  \item $\compdcm{B}{L}$, the computed (i.e., apparent) version of $\dcm{B}{L}$ is related additively to the true version by the error $\delta \dcm{B}{L}$
\[\compdcm{B}{L}= \dcm{B}{L}+ \delta \dcm{B}{L}\]
}
  \onslide<4->{
  \item It is also related to $\dcm{B}{L}$ by a small direction cosine matrix relating the local level frame to the computer frame $\widehat{L}$:
\[\compdcm{B}{L}= \dcm{L}{\widehat{L}}\dcm{B}{L}\] 
\end{itemize} 
}
\end{itemize}
\end{frame}

\begin{frame}
  \onslide<1->{
	\frametitle{Inertial Error Dynamics} 
	\framesubtitle{Strapdown Error Equations}
}
\begin{itemize}
  \onslide<2->{
\item Substituting the second equation into the first allows us to solve for $\delta \dcm{B}{L}$:
\begin{align*}
\delta \dcm{B}{L} &= \compdcm{B}{L} -\dcm{B}{L}\\
&= \dcm{L}{\widehat{L}}\dcm{B}{L}-\dcm{B}{L}\\
&= (\dcm{L}{\widehat{L}} - I)\dcm{B}{L}
\end{align*}
Now, let's introduce $\Phi=I-\dcm{L}{\widehat{L}}$ to obtain:
\[\delta \dcm{B}{L}=-\Phi\dcm{B}{L}\]
}
\end{itemize}
\end{frame}


\begin{frame}
  \onslide<1->{
	\frametitle{Inertial Error Dynamics} 
	\framesubtitle{Strapdown Error Equations}
}
\begin{itemize}
  \onslide<2->{
  \item Using small angle approximations, $\dcm{L}{\widehat{L}}$ can be written in the form
\[\dcm{L}{\widehat{L}}=
\begin{bmatrix}
1&\phi_z&-\phi_y\\
-\phi_z&1&\phi_x\\
\phi_y&-\phi_x&1
\end{bmatrix}.
\]
Thus, $\Phi=I-\dcm{L}{\widehat{L}}$ is seen to be a skew-symmetric rotation matrix:
\[\Phi=
\begin{bmatrix}
0&-\phi_z&\phi_y\\
\phi_z&0&-\phi_x\\
-\phi_y&\phi_x&0
\end{bmatrix}.
\]
}
\end{itemize}
\end{frame}


\begin{frame}
  \onslide<1->{
	\frametitle{Inertial Error Dynamics} 
	\framesubtitle{Strapdown Error Equations}
}
\begin{itemize}
  \onslide<2->{
  \item $\Phi$ is therefore equal to the cross product operator associated with a vector $\phi^L = \begin{bmatrix}\tiltx^L & \tilty^L & \az^L \end{bmatrix}\trans$.  I.e.,
\[\Phi = \xop{\phi^L}\]
}
  \onslide<3->{
\item The three components of $\phi^L$ will become three states in our inertial error dynamics Kalman filter.
}
  \onslide<4->{
\item $\phi_x$ and $\phi_y$ are called tilts, while $\psi_z$ is called the azimuth error.
}
\end{itemize}
\end{frame}


\begin{frame}
  \onslide<1->{
	\frametitle{Inertial Error Dynamics} 
	\framesubtitle{Strapdown Error Equations}
}
\begin{itemize}
  \onslide<2->{
  \item By a similar derivation, we can introduce a small angular position angle error $e^L=\begin{bmatrix}e_x^L & e_y^L & e_z^L \end{bmatrix}\trans$
}
  \onslide<3->{
  \item $e^L$ is associated with $\delta\dcm{E}{L}$ just as $\phi^L$ is associated with $\delta\dcm{B}{L}$:
\[\delta\dcm{E}{L}=-E^L\dcm{E}{L},\]
where $E^L$ is the skew-symmetric matrix form of $e^L$.
}
  \onslide<4->{
  \item The three components of $e^L$ yield the three position error states in a typical navigation Kalman filter.
}
\end{itemize}
\end{frame}


\begin{frame}
  \onslide<1->{
	\frametitle{Inertial Error Dynamics} 
	\framesubtitle{Strapdown Error Equations}
}
\begin{itemize}
  \onslide<2->{
  \item We are working to simplify Eq. \eqref{cblerreq}:
\[\cblerreq\]
}
  \onslide<3->{
  \item We have an expression for $\delta \dcm{B}{L}$.
}
\end{itemize}
\end{frame}


\begin{frame}
  \onslide<1->{
	\frametitle{Inertial Error Dynamics} 
	\framesubtitle{Strapdown Error Equations}
}
\begin{itemize}
  \onslide<2->{
  \item Consider the earth rate error term $\delta \rotmat{I}{E}^L$.  Since $\rotmat{I}{E}^L=\dcm{E}{L}\rotmat{I}{E}^E$ and $\delta\rotmat{I}{E}^E=0$, we derive:
\[\delta\rotmat{I}{E}^L=\delta\dcm{E}{L}\rotmat{I}{E}^E.\]
}
  \onslide<3->{
  \item But $\delta\dcm{E}{L}=-E^L\dcm{E}{L}$.  Therefore,
\begin{align*}
\delta\rotmat{I}{E}^L&=-E^L\dcm{E}{L}\rotmat{I}{E}^E\\
&=-E^L\rotmat{I}{E}^L.
\end{align*}
}
\end{itemize}
\end{frame}


\begin{frame}
  \onslide<1->{
	\frametitle{Inertial Error Dynamics} 
	\framesubtitle{Strapdown Error Equations}
}
\begin{itemize}
  \onslide<2->{
  \item For the error in transport rate, $\delta\rotmat{E}{L}^L$, using the error operator on Eq. \eqref{trateeq}, a new error terms having to do with altitude error shows up:
\begin{align*}
\delta\rotvec{E}{L}^L&=\delta\left[\dfrac{1}{R}u_r^L\times v^L\right]\\
&=\dfrac{1}{R}u_r^L\times \delta v^L - \dfrac{\delta h}{R^2}\left[u_r^L\times v^L\right]
\end{align*}
}
  \onslide<3->{
\item Here, $R=R_e+h$ for a spherical earth of radius $R_e$.
}
\end{itemize}
\end{frame}


\begin{frame}
  \onslide<1->{
	\frametitle{Inertial Error Dynamics} 
	\framesubtitle{}
}
\begin{itemize}
  \onslide<2->{
  \item Eq. \eqref{cblerreq} can now be rewritten as:
\begin{equation}\label{tilteq}
\tilteq
\end{equation} 
where $\delta \rotmat{E}{L}^L$ is the skew-symmetric matrix form of 
\[\trateerreqsphere\]
}
\end{itemize}
\end{frame}


\begin{frame}
  \onslide<1->{
	\frametitle{Inertial Error Dynamics} 
	\framesubtitle{Strapdown Error Equations}
}
\begin{itemize}
  \onslide<2->{
  \item Similarly, the velocity equation can be rewritten:
\begin{align*}
\velerreq
\end{align*} 
while the equation for the position angle error $e^L$ itself is
\begin{equation}
\paeeq
\end{equation} 
}
\end{itemize}
\end{frame}




\begin{frame}
  \onslide<1->{
	\frametitle{Inertial Error Dynamics} 
	\framesubtitle{Strapdown Error Equations}
}
\begin{itemize}
  \onslide<2->{
\item   Further Simplifications Using Telescope Pointing Angle Error

}
  \onslide<3->{
  \item $\psi^L = \phi^L - e^L$
}
  \onslide<4->{
  \item Leads to:
}
  \onslide<5->{
\begin{equation}\label{deltardoteq}
\deltardoteq
\end{equation}
and
\begin{eqnarray}
\velerreqpsi
\end{eqnarray}  
}
\end{itemize}
\end{frame}



\begin{frame}
  \onslide<1->{
	\frametitle{Inertial Error Dynamics} 
	\framesubtitle{Strapdown Error Equations}
}
\begin{itemize}
  \onslide<2->{
  \item Finally, the equations for the telescope pointing angle errors are: 
}
  \onslide<3->{
\begin{equation}\label{telscopeeq}
\telscopeeq
\end{equation}
}
\end{itemize}
\end{frame}

\begin{frame}
  \onslide<1->{
	\frametitle{Inertial Error Dynamics} 
	\framesubtitle{Strapdown Error Equations}
}
\begin{itemize}
  \onslide<2->{
 \item Three error equations yield three vector quantities of
 interest
 \begin{itemize}
}
  \onslide<3->{
 \item position error vector (3) $\poserr$ derived from $\delta \dcm{L}{E}$
}
  \onslide<4->{
 \item velocity error vector (3) $\velerr^L$ (or $\delta v^L$)
}
  \onslide<5->{
 \item attitude error vector (3) $\tiltx,\tilty,\az$ from $\delta \dcm{B}{L}$
 \end{itemize}
}
\end{itemize}
\end{frame}

\begin{frame}
  \onslide<1->{
	\frametitle{Inertial Error Dynamics} 
	\framesubtitle{Strapdown Error Equations}
}
\begin{itemize}
  \onslide<2->{
  \item See Matlab / Simulink Models
	\begin{itemize}
}
  \onslide<3->{
	  \item Strapdown Navigation Equations: TBD
}
  \onslide<4->{
	  \item Strapdown Navigation Error Equations: TBD
    \end{itemize}
}
\end{itemize}
\end{frame}

